\section[Introduction]{Introduction}

\begin{frame}{Problem Statement}
	\begin{itemize}
		\item So far, we have seen how models can make predictions for inputs $X$. We have also looked at how models behave globally through various metrics, and how to locally decompose error on test data. 
		\pause
		\item But for a specific model prediction where the true result is unknown, what is the error incurred?
		\pause
		\item In all application domains, a prediction error carries a \textbf{cost}.
	\end{itemize}
\end{frame}

\begin{frame}{Problem Statement}
	\begin{block}{Problem Statement}
		All models are approximations. Essentially, all models are wrong, but some are useful. However, the approximate nature of the model should never be lost from sight. \\
		\vspace{0.5cm}
		\hfill
		\textit{--- George E. P. Box (2007)}
	\end{block}
	\vfill 
	\tiny \textcolor{gray}{Source: Response Surfaces, Mixtures, and Ridge Analyses - George E. P. Box, Norman R. Draper - 2007 - p.414, John Wiley \& Sons}
\end{frame}

\begin{frame}{Prediction Regions and Confidence Intervals}
	In this course, we will introduce the concepts of prediction regions and confidence intervals for both classification and regression problems.
	\begin{itemize}
		\item First, we will develop an intuition for these intervals before providing a formal definition.
		\pause
		\item Next, we will see that linear models (regression and classification) provide theoretical estimates for these regions and intervals.
		\pause
		\item Finally, we will see how these regions and intervals can be empirically established for a wide variety of models.
	\end{itemize}
\end{frame}